\documentclass{article}

\usepackage{amsmath,amssymb,amsthm}
\usepackage{fullpage}
\usepackage{mathtools}

\newcommand{\RR}{\mathbb{R}}
\newcommand{\CC}{\mathbb{C}}
\newcommand{\QQ}{\mathbb{Q}}
\newcommand{\NN}{\mathbb{N}}
\newcommand{\ZZ}{\mathbb{Z}}
\newcommand{\PP}{\mathcal{P}}
\DeclareMathOperator{\GL}{GL}
\DeclareMathOperator{\SL}{SL}
\DeclareMathOperator{\cis}{cis}

\theoremstyle{definition}
\newtheorem*{solution}{Solution}

\title{Math 430 -- Problem Set 5 Solutions}
\date{Due April 1, 2016}

\begin{document}
\maketitle


\begin{enumerate}
\item[13.2.] Find all of the abelian groups of order $200$ up to isomorphism.
\begin{solution} 
Every abelian group is a direct product of cyclic groups.  Using the fact that $\ZZ_{mn} \cong \ZZ_m \times \ZZ_n$ if and only if $\gcd(m,n) = 1$, the list of groups of order $200$ is determined by the factorization of $200$ into primes:
\begin{itemize}
\item $\ZZ_8 \times \ZZ_{25}$
\item $\ZZ_4 \times \ZZ_2 \times \ZZ_{25}$
\item $\ZZ_2 \times \ZZ_2 \times \ZZ_2 \times \ZZ_{25}$
\item $\ZZ_8 \times \ZZ_5 \times \ZZ_5$
\item $\ZZ_4 \times \ZZ_2 \times \ZZ_5 \times \ZZ_5$
\item $\ZZ_2 \times \ZZ_2 \times \ZZ_2 \times \ZZ_5 \times \ZZ_5$.
\end{itemize}
\end{solution}
\item[13.5.] Show that the infinite direct product $G = \ZZ_2 \times \ZZ_2 \times \cdots$ is not finitely generated.
\begin{solution}
Note that every element of $G$ has order $2$ and that $G$ is abelian.  The group generated by any finite set of $k$ elements thus has at order at most $2^k$, while $G$ has infinite order.  Thus $G$ cannot be finitely generated.
\end{solution}
\item[16.1.] Which of the following sets are rings with respect to the usual operations of addition and multiplication?  If the set is a ring, is it also a field?
\begin{enumerate}
\item[(a)] $7\ZZ$
\begin{solution}
The is a subring of $\ZZ$ and thus a ring:
\begin{itemize}
\item $(7n) + (7m) = 7(m+n)$ so it is closed under addition;
\item $(7n)(7m) = 7(7mn)$ so it is closed under multiplication;
\item $-(7n) = (-7)(n)$, so it is closed under negation.
\end{itemize}
\end{solution}
It is not a field since it does not have an identity.
\item[(b)] $\ZZ_{18}$
\begin{solution}
This is a ring: the operations of arithmetic modulo $18$ are well defined.  It is not a field, since $2 \cdot 9 = 0$ gives a pair of zero divisors.
\end{solution}
\item[(c)] $\QQ(\sqrt{2})$
\begin{solution}
This is a subfield of $\RR$ and thus a field (in addition to being a ring:
\begin{itemize}
\item $(a + b\sqrt{2}) + (c + d\sqrt{2}) = (a + c) + (b + d)\sqrt{2}$ so it is closed under addition;
\item $(a + b\sqrt{2})(c+d\sqrt{2}) = (ac + 2bd) + (ad + bc)\sqrt{2}$ so it is closed under multiplication;
\item $-(a + b\sqrt{2}) = (-a) + (-b)\sqrt{2}$ so it is closed under negation;
\item $(a + b\sqrt{2})^{-1} = \frac{a}{a^2 - 2b^2} - \frac{b}{a^2 - 2b^2}\sqrt{2}$ and $a^2 - 2b^2 \ne 0$ for $a,b \in \QQ$ (unless both are zero)
\end{itemize}
\end{solution}
\item[(f)] $R = \{a + b\sqrt[3]{3} : a, b \in \QQ\}$
\begin{solution}
This is not a ring since $\sqrt[3]{3} \cdot \sqrt[3]{3}$ is not in $R$.
\end{solution}
\item[(h)] $\QQ(\sqrt[3]{3})$
\begin{solution}
This is a subfield of $\RR$ and thus a field:
\begin{itemize}
\item $(a + b\sqrt[3]{3} + c\sqrt[3]{9}) + (d + e\sqrt[3]{3} + f\sqrt[3]{9}) = (a+d) + (b+e)\sqrt[3]{3} + (c+f)\sqrt[3]{9}$ so it is closed under addition;
\item $(a + b\sqrt[3]{3} + c\sqrt[3]{9})(d + e\sqrt[3]{3} + f\sqrt[3]{9}) = (ad + 3bf + 3ce) + (ae + bd + 3cf)\sqrt[3]{3} + (af + be + cd)\sqrt[3]{9}$ so it is closed under multiplication;
\item $-(a + b\sqrt[3]{3} + c\sqrt[3]{9}) = (-a) + (-b)\sqrt[3]{3} + (-c)\sqrt[3]{9}$ so it is closed under negation.
\item Closure under inverses takes a bit more work, since there is no single conjugate.  We give two arguments.
\begin{enumerate}
\item In our first approach, we show directly that each element has an inverse. Given $(a + b\sqrt[3]{3} + c\sqrt[3]{9})$ we prove that there is some $(d + e\sqrt[3]{3} + f\sqrt[3]{9})$ with $(a + b\sqrt[3]{3} + c\sqrt[3]{9})(d + e\sqrt[3]{3} + f\sqrt[3]{9}) = 1$ using the formula for multiplication above.  This requires solving
\begin{align*}
ad + 3ce + 3bf &= 1 \\
bd + ae + 3cf &= 0 \\
cd + be + af &= 0.
\end{align*}
This system will have a solution as long as the determinant
\[
\det \begin{pmatrix}
a & 3c & 3b \\
b & a & 3c \\
c & b & a
\end{pmatrix} = a^3 + 3b^3 + 9c^3 - 9abc
\]
is nonzero.  Multiplying all of $a, b$ and $c$ by a common rational value we may assume that they are all integers and share no common factor.  If $a^3 + 3b^3 + 9c^3 - 9abc = 0$ then $a^3$ is a multiple of $3$, and thus $a$ is by unique factorization (say $a = 3a'$). Then
\[
9a'^3 + b^3 + 3c^3 - 9a'bc = 0.
\]
Repeating this for $b$ and $c$, we see that all are divisible by $3$, contradicting the fact that we chose them to have no common factor.
\item The other approach uses the extended Euclidean algorithm for polynomials.  Since $\sqrt[3]{3}$ is irrational, $x^3 - 3$ is irreducible and thus $\gcd(x^3 - 3, a + bx + cx^2) = 1$.  Therefore, there exist $f(x), g(x) \in \QQ[x]$ with
\[
f(x)(x^3 - 3) + g(x)(a + bx + cx^2) = 1.
\]
Evaluating this equation at $x = \sqrt[3]{3}$ shows that $g(\sqrt[3]{3})$ is the inverse of $a + b\sqrt[3]{3} + c\sqrt[3]{9}$.
\end{enumerate}
\end{itemize}
\end{solution}
\end{enumerate}
\item[16.2.] Let $R$ be the ring of $2 \times 2$ matrices of the form
\[
\begin{pmatrix} a & b \\ 0 & 0 \end{pmatrix},
\]
where $a, b \in \RR$. Show that although $R$ is a ring that has no identity, we can find a subring $S$ of $R$ with an identity.
\begin{solution}
Suppose that $\begin{psmallmatrix} a & b \\ 0 & 0\end{psmallmatrix}$ is an identity for $R$.  Then
\begin{align*}
\begin{pmatrix} 1 & 1 \\ 0 & 0 \end{pmatrix} \begin{pmatrix} a & b \\ 0 & 0 \end{pmatrix} &= \begin{pmatrix} 1 & 1 \\ 0 & 0 \end{pmatrix} \\
\begin{pmatrix} a & b \\ 0 & 0 \end{pmatrix} &= \begin{pmatrix} 1 & 1 \\ 0 & 0 \end{pmatrix},
\end{align*}
so $a=1$ and $b=1$.  But $\begin{psmallmatrix} 1 & 1 \\ 0 & 0\end{psmallmatrix}$ is not an identity, since
\[
\begin{pmatrix} 1 & 0 \\ 0 & 0 \end{pmatrix} \begin{pmatrix} 1 & 1 \\ 0 & 0 \end{pmatrix} = \begin{pmatrix} 1 & 1 \\ 0 & 0 \end{pmatrix}.
\]
Thus $R$ has no identity.

Let $S$ be the subring of matrices of the form $\begin{psmallmatrix} a & 0 \\ 0 & 0 \end{psmallmatrix}$.  Then $\begin{psmallmatrix} 1 & 0 \\ 0 & 0 \end{psmallmatrix}$ is an identity for $S$, since
\begin{align*}
\begin{pmatrix} 1 & 0 \\ 0 & 0 \end{pmatrix} \begin{pmatrix} a & 0 \\ 0 & 0 \end{pmatrix} &= \begin{pmatrix} a & 0 \\ 0 & 0 \end{pmatrix}, \\
\begin{pmatrix} a & 0 \\ 0 & 0 \end{pmatrix} \begin{pmatrix} 1 & 0 \\ 0 & 0 \end{pmatrix} &= \begin{pmatrix} a & 0 \\ 0 & 0 \end{pmatrix}.
\end{align*}
\end{solution}
\item[16.6.] Find all homomorphisms $\phi : \ZZ/6\ZZ \to \ZZ/15\ZZ$.
\begin{solution}
Since $\phi$ is a ring homomorphism, it must also be a group homomorphism (of additive groups).  Thuso $6\phi(1) = \phi(0) = 0$, and therefore $\phi(1) = 0, 5$ or $10$ (and $\phi$ is determined by $\phi(1)$).  If $\phi(1) = 5$, then
\begin{align*}
\phi(1) &= \phi(1 \cdot 1) \\
&= \phi(1) \cdot \phi(1) \\
&= 5 \cdot 5 \\
&= 10,
\end{align*}
which is a contradiction.  So the only two possibilities are
\begin{align*}
\phi_1(n) = 0 \mbox{ for all $n \in \ZZ_6$.}\\
\phi_2(n) = 10n \mbox{ for all $n \in \ZZ_6$.}
\end{align*}
We show that these are both well defined ring homomorphisms.

In both cases, adding a multiple of $6$ to $n$ changes the result by a multiple of $15$ ($0$ in the first case and $60$ in the second), so they are well defined.  They are additive group homomorphisms by the distributive law in $\ZZ_{15}$.  They are multiplicative since
\begin{align*}
0 \cdot 0 &= 0 \\
(10n) \cdot (10m) &= 100nm = 10nm.
\end{align*}
\end{solution}
\item[16.10.] Define a map $\phi : \CC \to \mathbb{M}_2(\RR)$ by
\[
\phi(a + bi) = \begin{pmatrix} a & b \\ -b & a \end{pmatrix}.
\]
Show that $\phi$ is an isomorphism of $\CC$ with its image in $\mathbb{M}_2(\RR)$.
\begin{solution}
We first show that $\phi$ is a homomorphism.
\begin{align*}
\phi((a + bi) + (c + di)) &= \phi((a + c) + (b + d)i) \\
&= \begin{pmatrix} a + c & b + d \\ -b - d & a + c \end{pmatrix} \\
\phi(a + bi) + \phi(c + di) &= \begin{pmatrix} a & b \\ -b & a \end{pmatrix} + \begin{pmatrix} c & d \\ -d & c \end{pmatrix} \\
&= \begin{pmatrix} a + c & b + d \\ -b - d & a + c \end{pmatrix} \\
\phi((a + bi)(c + di)) &= \phi((ac - bd) + (ad + bc)i) \\
&= \begin{pmatrix} ac - bd & ad + bc \\ -ad - bc & ac - bd \end{pmatrix} \\
\phi(a + bi)\phi(c + di) &= \begin{pmatrix} a & b \\ -b & a \end{pmatrix} \begin{pmatrix} c & d \\ -d & c \end{pmatrix} \\
&= \begin{pmatrix} ac - bd & ad + bc \\ -ad - bc & ac - bd \end{pmatrix}
\end{align*}
Finally, we show that $\phi$ is injective and thus an isomorphism onto its image.  If $\phi(a + bi) = \begin{psmallmatrix} 0 & 0 \\ 0 & 0 \end{psmallmatrix}$ then $a = 0$ and $b = 0$.
\end{solution}
\item[16.17.] Let $a$ be any element in a ring $R$ with identity.  Show that $(-1)a = -a$.
\begin{solution}
We have
\begin{align*}
(-1)a + a &= (-1)a + (1)a \\
&= (-1 + 1)a \\
&= (0)a \\
&= 0.
\end{align*}
The result now follows from the uniqueness of additive inverses.
\end{solution}
\item[16.22.] Prove the Correspondence Theorem: Let $I$ be an ideal of a ring $R$. Then $S \to S/I$ is a one-to-one correspondence between the set of subrings $S$ containing $I$ and the set of subrings of $R/I$. Furthermore, the ideals of $R$ correspond to the ideals of $R/I$.
\begin{solution} $ $
\begin{itemize}
\item We first show that the function $S \mapsto S/I$ sends subrings of $R$ to subrings of $R/I$.  If $s, t \in S$ then $(s + I) + (t + I) = (s + t) + I \in S/I$ since $S$ is closed under addition, $(s + I)(t + I) = (st) + I \in S/I$ since $S$ is closed under multiplication and $-(s + I) = (-s) + I \in S/I$ since $S$ is closed under negation.  Thus $S/I$ is a subring.
\item We now show that this function is surjective.  Let $T \subseteq R/I$ be a subring and set $S = \{x \in R : x + I \in T\}$.  Then $S$ is a subring of $R$: if $s, t \in S$ then $s+t \in S$ since $(s + t) + I = (s + I) + (t + I)$ and $T$ is closed under addition, $st \in S$ since $(st) + I = (s + I)(t + I)$ and $T$ is closed under multiplication, and $-s \in S$ since $(-s) + I = -(s + I)$ and $T$ is closed under negation.  Moreover, $S$ contains $I$ since $i + I = 0 + I \in T$ for every $i \in I$.  Finally, $S/I = T$ by construction.
\item Next, we show that this function is injective.  Suppose $S_1$ and $S_2$ are two subrings of $R$ that contain $I$ and that $S_1/I = S_2/I$ inside $R/I$.  We show that $S_1 \subseteq S_2$ (the opposite inclusion is analogous).  Suppose $x \in S_1$.  Since $S_1/I = S_2/I$, there is some $y \in S_2$ with $x + I = y + I$.  Thus there is an $i \in I$ with $x = y + i$.  Since $I \subseteq S_2$, both $y$ and $i$ are in $S_2$ and thus $x$ is as well.
\item Next, suppose that $S$ is actually an ideal of $R$.  To show that $S/I$ is an ideal of $R/I$, we take an arbitrary $s + I \in S/I$ and $x + I \in R/I$.  Then $xs + I \in S/I$ and $sx + I \in S/I$ since $S$ is an ideal of $R$.
\item Finally, suppose that $T \subseteq R/I$ is an ideal.  Let $S = \{x \in R : x + I \in T\}$ as before.  We show that $S$ is an ideal of $R$, proving that the correspondence restricts to a correspondence on ideals.  If $s \in S$ and $x \in R$ then $xs + I = (x+I)(s+I) \in S/I$ since $S/I$ is an ideal; similarly for $sx$.  Thus $S$ is an ideal.
\end{itemize}
\end{solution}
\item[16.26.] Let $R$ be an integral domain. Show that if the only ideals in $R$ are $\{0\}$ and $R$ itself, $R$ must be a field.
\begin{solution}
In order to show that $R$ is a field it suffices to prove that every nonzero $a \in R$ has an inverse.  Let $a \in R$ be nonzero, and consider the ideal $\langle a \rangle$ generated by $a$.  Since $a \ne 0$, this ideal is nonzero and thus is all of $R$ by assumption.  Therefore it contains $1$; by the definition of a principal ideal there is some $b \in R$ with $ab = 1$, providing the inverse of $a$.
\end{solution}
\item[16.31.] Let $R$ be a ring such that $1 = 0$.  Prove that $R = \{0\}$.
\begin{solution}
If $a \in R$ then $a = (1)a = (0)a = 0$, so $R = \{0\}$.
\end{solution}
\item[16.34.] Let $p$ be a prime. Prove that
\[
\ZZ_{(p)} = \{a/b : a, b \in \ZZ \mbox{ and $\gcd(b,p) = 1$}\}
\]
is a ring.
\begin{solution}
We show that this is a subring of $\QQ$ and thus a ring.  If $a/b, c/d \in \ZZ_{(p)}$ to show that $a/b + c/d = (ad + bc)/(bd)$ and $(a/b)(c/d) = ac)/(bd) \in \ZZ_{(p)}$ it suffices to show that $\gcd(bd, p) = 1$.  This follows from the fact that $p$ is prime: if it does not divide $b$ or $d$ then it cannot divide $bd$.  Negation is even easier: $-(a/b) = (-a)/b \in \ZZ_{(p)}$ since it has the same denominator.
\end{solution}
\item[16.38.] An element $x$ in a ring is called an idempotent if $x^2 = x$. Prove that the only idempotents in an integral domain are $0$ and $1$. Find a ring with an idempotent $x$ not equal to $0$ or $1$. 
\begin{solution}
If $x^2 = x$ then $(x-1)x = x^2 - x = 0$.  An integral domain has no zero divisors, so the only possibilities for $x$ are $0$ and $1$.

$3 \in \ZZ_6$ is an example of an idempotent that is neither $0$ nor $1$.
\end{solution}
\end{enumerate}

\end{document}